\subsection*{Part 5, Mission ``Say hello'': challenges}
The age question. \py|str( )| turns the number back into writing so it can be joined:
\begin{codenonum}
name = input("What's your name? ")
age = int(input("How old are you? "))
print("In 10 years you will be " + str(age + 10))
\end{codenonum}
\textbf{Super challenge:} without \py|int( )|, \py|age + 10| tries to add a number to writing. The last line says \texttt{TypeError: can only concatenate str (not "int") to str}. A \texttt{TypeError} means writing and numbers got mixed up.

\subsection*{Part 5, Mission ``Guess my number'': only 5 tries}
\begin{codenonum}
import random
secret = random.randint(1, 10)
guess = 0
tries = 0
while guess != secret and tries < 5:
    guess = int(input("Guess 1 to 10: "))
    tries += 1
    if guess < secret:
        print("Too low")
    elif guess > secret:
        print("Too high")
if guess == secret:
    print("You got it in " + str(tries) + " tries!")
else:
    print("Out of tries! It was " + str(secret))
\end{codenonum}
Typing \texttt{cat} gives \texttt{ValueError: invalid literal for int() with base 10: 'cat'}, because \py|int( )| can't turn the word \texttt{cat} into a number.

\subsection*{Part 5, Mission ``Quiz with a score'': the final message}
Put this at the very end, after the score is printed:
\begin{codenonum}
if score == 3:
    print("Quiz champion!")
elif score >= 1:
    print("Well done!")
else:
    print("Better luck next time!")
\end{codenonum}
\textbf{Super challenge:} put \py|import time| and \py|start = time.time()| at the top, and \py|print(round(time.time() - start))| at the bottom.

\subsection*{Part 5, Mission ``Reaction timer'': best of three}
\begin{codenonum}
import time
import random
best = 100
for i in range(3):
    input("Press Enter to start ")
    print("Get ready...")
    time.sleep(random.randint(2, 5))
    start = time.time()
    input("GO! Press Enter! ")
    taken = time.time() - start
    if taken < best:
        best = taken
print("Your best: " + str(round(best, 2)) + " seconds")
\end{codenonum}
\py|best| starts at 100 seconds, so your first go is sure to beat it.

\subsection*{Part 5: Translate it yourself}
Your code may look a little different and still be right. If it does the same thing, you did it!
\begin{codenonum}
answer = input("What's the password? ")
while answer != "banana":
    print("Wrong! Try again.")
    answer = input("What's the password? ")
print("Come in!")
\end{codenonum}
Using \py|while not answer == "banana":| on line 2 is right too.

\subsection*{Part 5: Two-line puzzle}
\begin{codenonum}
age = int(input("How old are you? "))
print("Next year you will be " + str(age + 1))
\end{codenonum}
